% !TeX root = ../Book1.tex
\chapter*{后记}
\addcontentsline{toc}{chapter}{后记}
回看本卷，最初的坐标变换已经有了几种不同的作用。在仿射与射影几何中，它帮助我们把对象化成标准形；在曲面和流形上，它把同一对象的局部表达联系起来；在向量丛与主丛中，它记录不同纤维标架怎样粘合。相似的矩阵公式由此承担了不同的几何任务。能够辨认公式正在改变什么、保留什么，是读完这些内容后值得保留的一种习惯。

\ProseParagraphBreak

局部计算与整体构造之间的距离，也在各编中逐渐显露。曲线的曲率、挠率等微分数据可以用来解出局部标架，曲线是否闭合还要检查走完一周后的结果；曲面的基本形式满足相容条件后，可以恢复局部浸入，沿不同道路延拓时又会遇到整体问题。到了平坦联络，曲率为零与环路平行移动非平凡可以同时发生。基本群和覆盖因此再次出现：早先用来辨别空间的工具，现在也记录几何数据的延拓。

\ProseParagraphBreak

本卷保留了古典几何、曲面和一般流形各自的篇幅。前者提供能够画出、计算和分类的对象，后者让这些经验适用于更广的空间。两者之间的联系有时要到第二次阅读才更清楚。例如，重新计算球面的曲率，可以同时观察局部坐标、移动标架与整体积分；重新看射影空间，可以把齐次坐标、商构造和群作用放在一起。读者也可以选定一个熟悉的空间，逐章积累它的不同描述，再尝试说明这些描述为何相容。

\ProseParagraphBreak

后续学习将沿本卷留下的具体问题展开。覆盖与曲面分类引出了更高维空间的不变量，第二卷用同调、上同调和同伦研究它们，并用横截性与临界点恢复光滑几何中的信息。度量、测地线与曲率引出了距离的整体比较和几何方程，第三卷继续处理这些问题。几乎复结构、微分形式和主丛中的相容条件，在第四卷的复几何、辛几何中有了新的内容；射影构造和由方程定义的几何对象，则在第五卷进入代数几何的讨论。这些方向各有进一步的基础要求，可以按所关心的问题接续阅读。

\ProseParagraphBreak

若希望继续研究本卷的专题，可以沿各章已给出的文献寻找更系统的处理：由曲面走向整体微分几何，由紧群走向表示论，由主丛与联络走向几何结构与规范理论。附录则保留了几条较短的延伸，例如用反射生成对称群、用有限构型研究多面体、用计算检验坐标公式。再次阅读时，选择一项能够独立完成的工作很有用：重建一个标准模型，算出一组过渡函数，或把一个局部解延拓到整个空间。遇到不能继续的地方，检查阻碍来自条件、拓扑还是计算，便能把下一步要学的问题说得更清楚。
